\documentclass[aspectratio=169,11pt]{beamer}
% Shared Beamer style for practice contest solution walkthroughs.
\usepackage[utf8]{inputenc}
\usepackage[T1]{fontenc}
\usepackage{lmodern}
\usepackage{amsmath,amssymb}
\usepackage{xcolor}
\usepackage{hyperref}
\usepackage{listings}

\definecolor{CPNavy}{HTML}{172033}
\definecolor{CPBlue}{HTML}{2563EB}
\definecolor{CPGreen}{HTML}{16A34A}
\definecolor{CPOrange}{HTML}{F97316}
\definecolor{CPGray}{HTML}{64748B}
\definecolor{CPLight}{HTML}{F8FAFC}
\definecolor{CPCodeBg}{HTML}{F1F5F9}
\definecolor{CPCodeRule}{HTML}{CBD5E1}

\usetheme{default}
\usefonttheme{professionalfonts}
\setbeamertemplate{navigation symbols}{}
\setbeamersize{text margin left=0.65cm,text margin right=0.65cm}
\setbeamertemplate{itemize item}{\color{CPBlue}$\blacktriangleright$}
\setbeamertemplate{itemize subitem}{\color{CPGray}$\triangleright$}

\setbeamercolor{normal text}{fg=CPNavy,bg=white}
\setbeamercolor{frametitle}{fg=white,bg=CPNavy}
\setbeamercolor{title}{fg=white,bg=CPNavy}
\setbeamercolor{subtitle}{fg=CPLight,bg=CPNavy}
\hypersetup{colorlinks=true,urlcolor=CPBlue}

\setbeamertemplate{frametitle}{%
  \nointerlineskip%
  \begin{beamercolorbox}[wd=\paperwidth,ht=0.88cm,dp=0.22cm,leftskip=0.55cm,rightskip=0.35cm]{frametitle}
    \usebeamerfont{frametitle}\insertframetitle\hfill\small\insertframenumber/\inserttotalframenumber
  \end{beamercolorbox}%
}

\setbeamertemplate{footline}{%
  \leavevmode%
  \hbox{%
    \begin{beamercolorbox}[wd=.58\paperwidth,ht=2.4ex,dp=1ex,leftskip=.5cm]{author in head/foot}%
      \color{CPGray}\insertshorttitle
    \end{beamercolorbox}%
    \begin{beamercolorbox}[wd=.42\paperwidth,ht=2.4ex,dp=1ex,rightskip=.5cm plus1fil]{date in head/foot}%
      \hfill\color{CPGray}\insertshortauthor
    \end{beamercolorbox}%
  }%
}

\setbeamertemplate{title page}{%
  \vspace*{1.25cm}
  \begin{beamercolorbox}[wd=\paperwidth,sep=0.65cm,leftskip=0.15cm,rightskip=0.15cm]{title}
    {\color{white}\Huge\bfseries\inserttitle\par}
    \vspace{0.35cm}
    {\color{CPLight}\Large\insertsubtitle\par}
    \vspace{0.6cm}
    {\color{CPLight}\large\insertauthor\par}
  \end{beamercolorbox}
  \vfill
}

\newcommand{\sectiontag}[1]{\textbf{\color{CPBlue}#1}}
\newenvironment{tightitemize}{\begin{itemize}\small\setlength{\itemsep}{0.18cm}}{\end{itemize}}
\lstdefinestyle{cpcode}{
  language=Python,
  basicstyle=\ttfamily\tiny,
  keywordstyle=\color{CPBlue}\bfseries,
  commentstyle=\color{CPGray}\itshape,
  stringstyle=\color{CPGreen},
  backgroundcolor=\color{CPCodeBg},
  frame=single,
  framerule=0.25pt,
  rulecolor=\color{CPCodeRule},
  showstringspaces=false,
  columns=fullflexible,
  keepspaces=true,
  breaklines=true,
  breakatwhitespace=false,
  tabsize=4,
  xleftmargin=0.08cm,
  xrightmargin=0.08cm,
  aboveskip=0.08cm,
  belowskip=0.08cm
}


\title[firefly]{Firefly}
\subtitle{Day 2 solution walkthrough}
\author[Solution Deck]{Competitive Programming Bootcamp}
\date{}

\begin{document}

\begin{frame}[plain]
  \titlepage
\end{frame}

% BEGIN GENERATED STATEMENT INTERPRETATION
\begin{frame}{Interpret the Word Problem}
\small
\sectiontag{Statement excerpts:}
\begin{tightitemize}
  \item \textit{``choose a single height''}
  \item \textit{``minimum number of obstacles the firefly needs to destroy''}
\end{tightitemize}

\vspace{0.18cm}
\sectiontag{Programming interpretation:}
\begin{tightitemize}
  \item Each possible height can be evaluated independently.
  \item Bottom and top obstacles use different hit thresholds.
  \item Sorting both obstacle lists lets each threshold count be answered by binary search.
\end{tightitemize}
\end{frame}
% END GENERATED STATEMENT INTERPRETATION







\begin{frame}{Problem Model}
\small
\sectiontag{Textbook connection:} CP4 3.3.1 Binary Search

\vspace{0.25cm}
\sectiontag{Core technique:} Sorted arrays plus binary search

\vspace{0.25cm}
\begin{tightitemize}
  \item Obstacles alternate between bottom stalagmites and top stalactites.
  \item For every possible flying height, count how many obstacles are hit.
  \item Return the minimum hit count and how many heights achieve it.
\end{tightitemize}
\end{frame}

\begin{frame}{How to Recognize the Approach}
\begin{tightitemize}
  \item For a fixed height, an obstacle hits if its length is at least a threshold.
  \item Counting values above a threshold in a sorted list is a binary search problem.
  \item Bottom and top obstacles have different threshold formulas.
\end{tightitemize}
\end{frame}

\begin{frame}{Algorithm}
\begin{tightitemize}
  \item Split obstacle lengths into bottom and top lists by parity of input index.
  \item Sort both lists.
  \item For each height h, bottom hits are lengths >= h.
  \item Top hits are lengths >= H - h + 1.
  \item Use bisect\_left to count each group and update the minimum/count pair.
\end{tightitemize}
\end{frame}

\begin{frame}{Walkthrough of the Key State}
\begin{tightitemize}
  \item bisect\_left(list, threshold) returns the first index that can hit.
  \item len(list) - index is the number of hitting obstacles.
  \item Trying all heights is safe because each height query is logarithmic.
\end{tightitemize}
\end{frame}

\begin{frame}{Why the Algorithm Is Correct}
\begin{tightitemize}
  \item Every possible height is evaluated exactly once.
  \item The binary searches compute exact hit counts because both lists are sorted.
  \item The maintained minimum and frequency therefore describe the best heights over the full range.
\end{tightitemize}
\end{frame}

\begin{frame}{Complexity and Implementation Notes}
\small
\sectiontag{Complexity:} O(N log N + H log N) time and O(N) memory.

\vspace{0.3cm}
\begin{tightitemize}
  \item The top threshold uses H - h + 1 because heights are 1-indexed from the floor.
  \item Keep the two obstacle types separate; a single sorted list loses direction information.
\end{tightitemize}
\end{frame}



% BEGIN GENERATED CODE WALKTHROUGH
\begin{frame}[fragile]{Code: Read Cave Dimensions}
\scriptsize
\sectiontag{Source lines:} \texttt{firefly.py:41--52}

\vspace{0.08cm}
\begin{columns}[T,totalwidth=\textwidth]
\begin{column}{0.58\textwidth}
\begin{lstlisting}[style=cpcode]
import sys
from bisect import bisect_left

def main():
    input = sys.stdin.buffer.readline

    n, h = map(int, input().split())

    bottom = []
    top = []
\end{lstlisting}
\end{column}
\begin{column}{0.38\textwidth}
\sectiontag{What this chunk does}
\begin{tightitemize}
  \item n is the number of obstacles and h is the cave height.
  \item Two lists separate obstacles that grow from the floor and hang from the ceiling.
  \item This separation matches the two different hit formulas.
\end{tightitemize}
\end{column}
\end{columns}
\end{frame}

\begin{frame}[fragile]{Code: Split and Sort Obstacles}
\scriptsize
\sectiontag{Source lines:} \texttt{firefly.py:53--64}

\vspace{0.08cm}
\begin{columns}[T,totalwidth=\textwidth]
\begin{column}{0.58\textwidth}
\begin{lstlisting}[style=cpcode]
for i in range(n):
    length = int(input())

    if i % 2 == 0:
        bottom.append(length)
    else:
        top.append(length)

bottom.sort()
top.sort()
\end{lstlisting}
\end{column}
\begin{column}{0.38\textwidth}
\sectiontag{What this chunk does}
\begin{tightitemize}
  \item The problem states obstacles alternate bottom, top, bottom, top.
  \item Even positions enter bottom; odd positions enter top.
  \item Sorting is required before bisect\_left can count obstacles by threshold.
\end{tightitemize}
\end{column}
\end{columns}
\end{frame}

\begin{frame}[fragile]{Code: Count Hits by Height}
\scriptsize
\sectiontag{Source lines:} \texttt{firefly.py:66--82}

\vspace{0.08cm}
\begin{columns}[T,totalwidth=\textwidth]
\begin{column}{0.58\textwidth}
\begin{lstlisting}[style=cpcode]
minimum = n
count = 0

for height in range(1, h + 1):

    index = bisect_left(bottom, height)
    bottom_hits = len(bottom) - index

    required_length = h - height + 1

    index = bisect_left(top, required_length)
    top_hits = len(top) - index

    hits = bottom_hits + top_hits
\end{lstlisting}
\end{column}
\begin{column}{0.38\textwidth}
\sectiontag{What this chunk does}
\begin{tightitemize}
  \item For a bottom obstacle, length at least height means it is hit.
  \item For a top obstacle, length at least h - height + 1 means it reaches down to that level.
  \item len(list) - index gives the number of obstacles meeting the threshold.
\end{tightitemize}
\end{column}
\end{columns}
\end{frame}

\begin{frame}[fragile]{Code: Track the Best Levels}
\scriptsize
\sectiontag{Source lines:} \texttt{firefly.py:84--94}

\vspace{0.08cm}
\begin{columns}[T,totalwidth=\textwidth]
\begin{column}{0.58\textwidth}
\begin{lstlisting}[style=cpcode]
        if hits < minimum:
            minimum = hits
            count = 1

        elif hits == minimum:
            count += 1

    print(minimum, count)

main()
\end{lstlisting}
\end{column}
\begin{column}{0.38\textwidth}
\sectiontag{What this chunk does}
\begin{tightitemize}
  \item A lower hit count resets both the minimum and the number of best heights.
  \item A tie increments the count of optimal heights.
  \item The final output is the minimum obstacles and the number of heights achieving it.
\end{tightitemize}
\end{column}
\end{columns}
\end{frame}
% END GENERATED CODE WALKTHROUGH

\begin{frame}{Source}
\small
\sectiontag{Solution file:} \texttt{practice\_contests/day\_2/firefly.py}

\vspace{0.3cm}
\sectiontag{Problem:} \url{https://open.kattis.com/problems/firefly}

\vspace{0.45cm}
Use this deck with the source code open beside it. The slides explain the reasoning; the code shows the exact input handling and output formatting.
\end{frame}

\end{document}
